% Frozen after the additive projection failed on old counts, before new Aer.
\documentclass[11pt,a4paper]{article}
\usepackage{kotex}
\usepackage{amsmath,amssymb,booktabs,geometry}
\geometry{margin=25mm}
\title{I01 선검증 4b 사전 정정: 경계 보존 합산-null bootstrap}
\author{RELATIONAL ATLAS}
\date{2026년 10월 4일}
\begin{document}
\maketitle

\section{수정 경위}
선검증 4의 가법적 null 투영을 기존 raw counts에 적용한 결과,
일부 표지의 추정 성공확률이 $[0,1]$ 밖으로 나갔다.
원래 프로토콜의 중단 규칙에 따라 새 Aer seed는 실행하지 않았다.
구 프로토콜과 실패 코드는 남겨 두고, 아래의 경계 보존 투영을
\emph{새 Aer 실행 전에} 고정한다. 기존 raw counts의 재분석은
사후 진단으로만 보고한다.

\section{회로·표본·검정 가족}
회로, $n=3$, $L=7$, $\beta=0.45$, $\beta_r=0.15$,
$E=E_0+0.37E_0^2/3$, $\tau=0,\pi/2$, 공통 위상
$0\to0.60$ rad, 균일/비균일 표지 배분, 비제약/$P=0$ 상태,
표지별 shots 및 총 336회로는 선검증 3과 같다. CPU workers 6,
독립 Aer seed 20261020이다. 가족은 여덟 약한 합산 null 검정이고
Bonferroni 문턱은 $0.00625$다. 일차 통계량은 선검증 4의
표지 가중 차이 $\hat\Delta_i$, 이항 분산 $\hat V_i$ 및
$T=\sum_i\hat\Delta_i^2/\hat V_i$를 그대로 사용한다.

\section{경계 보존 null 생성}
축·표지·블록의 성공확률 초기 추정은 Jeffreys 보정
$r_{b,it}=(c_{b,it}+1/2)/(N_t+1)$을 쓴다.
$\operatorname{logit}r=\log[r/(1-r)]$로 두고 각 축에 공통
실수 $u_i$를 찾아
\begin{equation}
\sum_t a_t\left[
\operatorname{logistic}(\operatorname{logit}r_{1,it}-u_i)-
\operatorname{logistic}(\operatorname{logit}r_{0,it}+u_i)
\right]=0
\end{equation}
을 만족시킨다. 이 식의 좌변은 $u_i$에 대해 단조 감소하고
양 끝에서 부호가 달라 이분법으로 푼다.
그 결과를 $p'_{0,it},p'_{1,it}$로 사용하면 모든 확률은 $(0,1)$에
있고 가중 합 차이는 수치 오차에서 0이다.
독립 이항표본 10000패널마다 $T^*$를 다시 계산하고
$(1+\#\{T^*\ge T\})/10001$을 주 $p$값으로 쓴다.
$\chi^2_3$ 꼬리확률은 별도 교차 점검으로 기록한다.
기존 raw counts의 사후 재분석과 새 seed의 사전 검정 결과를
서로 섞지 않는다. QPU 및 발전연구 주제 선정은 보류한다.
\end{document}
